% =============================================================================
% Automated Cost-Surface-Driven Network Routing for Rural Cluster Connectivity
% and Downstream Impact Analysis of Building Detection Omissions
% Compile: pdflatex main && bibtex main && pdflatex main && pdflatex main
% All numerical results in Sections 5.2-5.4 are taken from the project outputs
% in DELIVERABLES_DAMARAKUNTA_STAGE_B/ and docs/results/ (see docs/results/PROVENANCE_AND_PROOFS.md).
% =============================================================================
\documentclass[journal]{IEEEtran}

\usepackage{amsmath,amssymb,amsfonts}
\usepackage{graphicx}
\usepackage{booktabs}
\usepackage{array}
\usepackage{multirow}
\usepackage{siunitx}
\usepackage[ruled,vlined,linesnumbered]{algorithm2e}
\usepackage{cite}
\usepackage{textcomp}
\usepackage[hidelinks]{hyperref}

\newcommand{\Lcost}{L_{\mathrm{cost}}}
\newcommand{\dinoseg}{StrongDINOv2Seg}
\DeclareMathOperator*{\argmin}{arg\,min}
\DeclareMathOperator*{\argmax}{arg\,max}
\sisetup{group-separator={,}, group-minimum-digits=4}

\begin{document}

\title{Automated Cost-Surface-Driven Network Routing for Rural Cluster Connectivity and Downstream Impact Analysis of Building Detection Omissions}

\author{Pendyala Archana%
\thanks{P.~Archana is with the Chaitanya Bharathi Institute of Technology (CBIT), Hyderabad, Telangana, India.}}

\markboth{Manuscript draft}{P. Archana: Cost-surface routing and building-omission impact analysis}
\maketitle

% -----------------------------------------------------------------------------
\begin{abstract}
\textbf{Background:} Habitation-level rural road programmes need building inventories that are still digitised by hand, and automated extraction is judged by pixel scores (IoU, Dice) that say little about planning impact; planners also need the same logic at Gram Panchayat, mandal and district scale. \textbf{Methods:} We couple a fine-tuned DINOv2 segmenter (StrongDINOv2Seg) with a 10~m multi-criteria cost surface (land cover, slope, footprint penalty, local penalty override) and a shared-network planner, apply it to Damarakunta (Siddipet district, Telangana), propagate controlled building omissions through the plan, and extend it hierarchically to Gram Panchayat, mandal and district tiers (coarse-to-fine routing with graph contraction). \textbf{Results:} StrongDINOv2Seg reaches 53.8\% mIoU, 68.9\% Dice and 71.3\% pixel accuracy on LoveDA validation. At village level the shared network connects 163 houses with 5{,}069~m of corridor, 24.9\% less than independent routes; omitting 10--20\% of candidates understates length by 7.6--16.9\% and leaves 16--33 households unplanned, and omitting 62 merged blobs leaves 38.0\% unserved. At Gram Panchayat level a footprint penalty of 4 removes 89.0\% of route length inside footprints for 5.2\% more length. Mandal and district runs on proxy nodes demonstrate scalability: the district backbone is 95.9\% existing road, with 14.7~km of new construction. \textbf{Conclusion:} Pixel accuracy is an insufficient proxy for planning utility; omission-aware, per-target reporting should accompany detection metrics. Village and Gram Panchayat results use real inputs; the mandal and district runs use proxy hubs and are not planning findings.
\end{abstract}

\begin{IEEEkeywords}
Building footprint extraction, vision foundation model, DINOv2, least-cost path, cost surface, Steiner tree, rural road planning, error propagation, PMGSY.
\end{IEEEkeywords}

% -----------------------------------------------------------------------------
\section{Introduction}
\IEEEPARstart{R}{ural} road connectivity remains a binding constraint on access to health care, schooling and markets across developing regions. In India, the Pradhan Mantri Gram Sadak Yojana (PMGSY) prioritises all-weather connection of unconnected habitations by population threshold~\cite{pmgsy}. Planning such connections at the level of individual hamlets requires (i) an inventory of dwellings and priority facilities, (ii) a model of the terrain and land cover that the corridor must cross, and (iii) a method for turning these into candidate alignments. In practice all three are assembled manually: dwellings are digitised from imagery, corridors are drawn by engineers, and each house is treated as a separate connection problem. Manual digitisation is slow and subjective, and independent house-by-house connections duplicate corridor length where a shared access road would serve several households.

Deep learning has made large-scale building extraction feasible~\cite{ji2019whu,sirko2021openbuildings}, and vision foundation models promise better transfer to sparse, heterogeneous rural scenes~\cite{oquab2023dinov2,cong2022satmae}. Yet segmentation research is evaluated almost exclusively with pixel-level overlap measures. A segmentation map is rarely the end product: here it is the input to a graph problem whose output is a road alignment and a length estimate. A building that is missed is not merely a few lost pixels; it may be an unconnected household, a shorter network, and a corridor that crosses a roof the planner did not know existed. The relation between IoU or Dice and such downstream quantities is neither monotone nor well characterised, and it is invisible when the two stages are reported separately~\cite{heuvelink1998error,foody2002status}.

This paper closes that loop for a concrete rural case. We build an end-to-end pipeline from satellite imagery to shared access roads for Damarakunta, Siddipet district, Telangana, and use controlled omission scenarios to quantify how detection false negatives propagate into topology, length and household coverage. The contributions are:
\begin{itemize}
\item \textbf{A foundation-model segmentation stage.} StrongDINOv2Seg, a fully fine-tuned DINOv2 encoder with a multi-layer FPN/ASPP decoder, trained with a composite loss and validated on LoveDA, whose epoch-4 checkpoint produced the building candidates used downstream.
\item \textbf{A multi-criteria cost surface with a local penalty override.} Land-cover resistance, slope limits and a soft building-footprint penalty are combined on a 10~m grid; a spatially bounded override lifts the footprint penalty around a priority facility whose point lies on the edge of its own footprint.
\item \textbf{A two-stage shared-network planner.} Stage~A generates up to $K$ distinct least-cost alternatives per target; Stage~B selects among them greedily by marginal cost with shared-edge credit, followed by leaf re-routing. We compare it with independent routing, a length-only control, a straight-line baseline and a roads-only shared reference.
\item \textbf{An omission-propagation analysis.} Deterministic and random omission and addition scenarios quantify changes in network length, length per routed target, unserved households and route length inside omitted footprints, separating the planner's information loss from structural restructuring using a no-penalty control.
\item \textbf{A hierarchical multi-scale extension.} The same cost-surface and shared-network logic is lifted to Gram Panchayat (10~m, local overrides), mandal (30~m, collector-road attachment) and district (250~m coarse plan, graph contraction, windowed 30~m refinement) tiers, together with a footprint-penalty sensitivity sweep ($\mu\in\{1,2,4,8,20,100\}$) that supports $\mu=4$. The mandal and district tiers are reported strictly as a computational scalability demonstration on proxy nodes.
\item \textbf{Planning guidance.} We argue for reporting omission-aware, per-target quantities alongside IoU and Dice, and discuss implications for PMGSY-style programmes.
\end{itemize}

The remainder of the paper is organised as follows. Section~\ref{sec:related} reviews related work, Section~\ref{sec:data} describes the study area and data, Section~\ref{sec:method} presents the methodology, Section~\ref{sec:results} reports results, and Section~\ref{sec:conclusion} concludes.

% -----------------------------------------------------------------------------
\section{Related Work}\label{sec:related}

\subsection{Deep Learning for Rural Building Extraction}
Encoder--decoder networks descending from U-Net~\cite{ronneberger2015unet} dominate building and road extraction. Variants add nested skip pathways~\cite{zhou2018unetpp}, residual units~\cite{zhang2018resunet}, atrous multi-scale context~\cite{chen2018deeplabv3plus} or feature pyramids~\cite{lin2017fpn}. Deep supervision~\cite{lee2015dsn} improves gradient flow and boundary quality in decoders. Benchmarks such as WHU~\cite{ji2019whu} and LoveDA~\cite{wang2021loveda} have been decisive; LoveDA is notable for its explicit urban/rural domain split. In earlier prototyping for this project, standard CNN architectures (U-Net, UNet++, PSPNet, LinkNet and DeepLabV3+ from the \texttt{segmentation\_models\_pytorch} library with ResNet-34 encoders, plus HRNet and FCN-8s) were set up as the architectural exploration context; StrongDINOv2Seg was selected for the final deployment on Damarakunta, and no results of those CNN baselines are reported here. Continental-scale products such as Open Buildings~\cite{sirko2021openbuildings} show that rural structures can be mapped at scale, but also that thatched, tin-roofed, densely clustered and tree-occluded dwellings remain the hardest cases, and that detectors disagree substantially where reference data are absent.

Vision foundation models change the trade-off. Vision transformers trained with self-supervision, notably DINOv2~\cite{oquab2023dinov2}, give dense patch features that transfer to segmentation with little task-specific data. Remote-sensing-specific pre-training, such as SatMAE~\cite{cong2022satmae} and Scale-MAE~\cite{reed2023scalemae}, adds temporal and scale awareness. The cost is a coarse patch grid (14~px for DINOv2) that blurs the small, closely spaced footprints typical of rural hamlets; decoders must recover boundaries, which motivates the multi-layer fusion and boundary-weighted loss used here.

\subsection{Automated Road Network Generation and Least-Cost Routing}
Least-cost path (LCP) analysis on a raster friction surface is the standard tool for alignment planning. Shortest-path search on grid graphs~\cite{hart1968astar} is ubiquitous, and direction-dependent formulations account for slope-sensitive travel~\cite{collischonn2000lcp}. Network libraries such as OSMnx~\cite{boeing2017osmnx,boeing2025osmnx2} make road graphs from OpenStreetMap directly usable for such analyses. Connecting \emph{several} terminals at minimum total length is the Steiner tree problem~\cite{hwang1992steiner}, for which metric-closure approximations~\cite{kou1981steiner,mehlhorn1988steiner} give constant-factor guarantees. Planning studies typically route each destination independently or solve a Steiner relaxation on a sparse graph; few operate on a dense, penalised cost raster where footprints, facilities and an existing road network interact, and few examine how input inventory errors affect the resulting alignments.

\subsection{Error Propagation in Geospatial Pipelines}
The propagation of input uncertainty through GIS operations has a long history~\cite{heuvelink1998error}, and map accuracy assessment is a field of its own~\cite{foody2002status}. Global measures such as overall accuracy do not predict the error of downstream aggregates, because errors are spatially structured and the downstream operator is nonlinear. In graph-based planning a missing node can disconnect a household, a missing obstacle can invite a route through a dwelling, and a missing cluster can understate the corridor budget. We found no study that couples a modern segmentation pipeline with a shared-network planner and measures these effects; this paper is a first step on one village.

% -----------------------------------------------------------------------------
\section{Study Area and Dataset}\label{sec:data}

\subsection{Damarakunta}
Damarakunta is a rural revenue village in Siddipet district, Telangana (approximately \ang{17.707}\,N, \ang{78.753}\,E). The administrative boundary, obtained from the Telangana State Remote Sensing Application Centre (TGRAC), encloses $18.47$~km$^2$ and is processed in UTM zone 44N (EPSG:32644). Settlement is dispersed: houses form small clusters, often along or off a few tracks, separated by agricultural land. The cadastral layer contains 794 parcels. Four facilities are recorded by TGRAC (two schools, a health sub-centre and an anganwadi); the anganwadi was excluded because its coordinate lies 371~m from the nearest detected footprint and 126~m from the nearest independent building polygon. Two Gram Panchayat headquarters points were added as administrative utilities, giving five facilities in the planning layer. The TGRAC road layer has 43 features totalling $93.5$~km, of which 44.7\% lies outside the village boundary.

\subsection{Raster and Vector Data}
\textbf{Imagery.} A $0.6$~m, three-band mosaic ($11{,}296 \times 7{,}171$ px, EPSG:32644) was assembled from high-resolution web basemap exports. \textbf{Terrain.} Slope is computed with a Sobel operator from a $\sim$30~m DEM and bilinearly resampled to 10~m; it therefore carries no genuine 10~m terrain detail. \textbf{Land cover.} ESA WorldCover~\cite{zanaga2021worldcover} supplies land-cover classes at 10~m. \textbf{Roads and cadastre.} TGRAC reference roads were clipped to the boundary (51.7~km retained); cadastral parcels were repaired with a validity operation and are used only to count parcel crossings, never as routing constraints. \textbf{Grid.} All layers are aligned to a $432 \times 679$ cell grid at 10~m. Cells that are water, no-data or steeper than $25^\circ$ are blocked; 961 slope and 178 land-cover cells lacking data are treated as invalid rather than interpolated.

\subsection{Building Candidates and Annotation Protocol}
The building candidates are derived from StrongDINOv2Seg predictions: the project's tiled-inference script applies the epoch-4 checkpoint (\texttt{best\_model.pth}) to the mosaic in $560\times560$ px tiles (stride 420 px, 70 px margin), and the raw class maps were vectorised and post-processed into instance-level \emph{building candidates}. The checkpoint file's hash was not logged with the candidate layer, so this provenance rests on the project scripts and on the checkpoint's stored metadata (epoch 4, mIoU 0.5380, mDice 0.6891), not on a recorded hash. The layer holds 278 polygons, 165 of them inside the boundary and 163 retained as routing targets. Each candidate carries a quality class: \textsc{Plausible} (71), \textsc{Review} (32) and \textsc{Unresolved\_Merge} (62) for the 165 inside the boundary, where the last class denotes blobs that merge several adjacent roofs into one polygon and could not be split automatically.

No pixel-level reference exists for Damarakunta: the project has neither a manual building mask nor a field survey. Two evaluations are therefore used and must be kept distinct. First, \emph{source-domain} validation of the segmentation model on LoveDA (Section~\ref{sec:segres}). Second, a \emph{proxy agreement} check of the village candidates against Google Open Buildings v3~\cite{sirko2021openbuildings}, a different detector on other imagery; a candidate agrees if an Open Buildings polygon lies within 30~m, and a facility is supported if a footprint lies within 20~m of its point. Agreement is not accuracy, and neither layer is treated as truth. A visual-audit protocol (imagery crops and a decision rule: confirm within 20~m, relocate within 200~m, otherwise reject) has been prepared for the facilities but has not yet been completed in the field or in imagery.

% -----------------------------------------------------------------------------
\subsection{Open Data for the Hierarchical Levels}\label{sec:opendata}
Levels 2--4 (Section~\ref{sec:hier}) use only open layers, fetched by a script that records the URL, release, licence and SHA-256 of every file. \textbf{Land cover:} ESA WorldCover 10~m v200 (2021)~\cite{zanaga2021worldcover}. \textbf{Terrain:} the Copernicus GLO-30 DEM~\cite{esa2022copdem}, read as windows and reprojected to EPSG:32644; slope is computed with the Stage~B method (Sobel on the 30~m grid, bilinear to the target grid). Stage~B itself used the SRTM 1~arc-second DEM, which is not in the repository; over the Damarakunta window the two DEMs differ by an RMSE of 2.3~m (bias $-1.6$~m, $r=0.996$), the $>25^\circ$ hard-block mask differs in no cell, and re-running Level~2 with either DEM changes the shared length by 1.0\% and the ordinal cost by 2.6\%, with an identical override detour reduction (the SRTM used for this comparison is a public SRTM-derived tile set, a proxy for the original file). \textbf{Boundaries:} geoBoundaries gbOpen ADM2/ADM3~\cite{runfola2020geoboundaries}, giving the Siddipet district polygon (3{,}722~km$^2$), its 23 mandal polygons (as released; the current administrative set may differ) and Markook mandal (116~km$^2$). \textbf{Roads and villages:} Overture Maps transportation segments and village localities~\cite{overture2026}, release 2026-10-21.0 (OSM-derived), restricted to collector-grade classes (motorway to tertiary) as attachment targets at Levels~3--4. Roads and villages from Overture are not a PMGSY or TGRAC inventory, and no open Gram Panchayat polygons exist. Hence the mandal and district tiers use \emph{proxy} nodes: village points as Gram Panchayat hubs (weight 1, house counts unknown), the interior point of the Markook polygon as the Mandal HQ, and the interior points of the 23 mandal polygons as terminals; none is a verified HQ, PHC or market location. Google Open Buildings v3~\cite{sirko2021openbuildings} is used only for the proxy agreement check of Section~\ref{sec:data}.

\section{Methodology}\label{sec:method}

\subsection{Pipeline Overview}
The pipeline has four stages: (1) building and facility extraction from imagery; (2) construction of a multi-criteria cost surface; (3) shared network generation by least-cost path alternatives and greedy selection; and (4) sensitivity analysis under controlled omissions and additions. Stages 2--3 consume the output of Stage~1 only through the footprint polygons and target points, which is what makes stage-wise error propagation measurable.

\subsection{Segmentation Model}

\paragraph{Problem and loss.}
Let $x\in\mathbb{R}^{H\times W\times 3}$ be an image tile and $y\in\{0,\dots,C-1\}^{H\times W}$ its label map, with class $0$ ignored and $C=8$ (LoveDA convention). For logits $z$ and softmax probabilities $p_c$, we minimise a composite loss
\begin{equation}
\mathcal{L}_{\mathrm{seg}} = \lambda_{\mathrm{ce}}\mathcal{L}_{\mathrm{ce}} + \lambda_{\mathrm{dice}}\mathcal{L}_{\mathrm{dice}} + \lambda_{\mathrm{edge}}\mathcal{L}_{\mathrm{edge}},
\label{eq:seg}
\end{equation}
with $(\lambda_{\mathrm{ce}},\lambda_{\mathrm{dice}},\lambda_{\mathrm{edge}})=(0.5,0.4,0.1)$. Cross-entropy is $\mathcal{L}_{\mathrm{ce}}=-\frac{1}{|\Omega|}\sum_{i\in\Omega}\log p_{y_i}(i)$ over non-ignored pixels $\Omega$. The soft Dice term~\cite{milletari2016vnet} is
\begin{equation}
\mathcal{L}_{\mathrm{dice}} = \frac{1}{C-1}\sum_{c=1}^{C-1}\Bigl(1-\frac{2\sum_i p_c(i)\,\mathbb{1}[y_i=c]+\epsilon}{\sum_i p_c(i)+\sum_i\mathbb{1}[y_i=c]+\epsilon}\Bigr),
\end{equation}
and the boundary term reweights cross-entropy near label edges,
\begin{equation}
\mathcal{L}_{\mathrm{edge}} = \frac{1}{|\Omega|}\sum_{i\in\Omega}\bigl(1+4\,e_i\bigr)\,\ell_{\mathrm{ce}}(i),\quad e_i=\min\bigl(1,\,|G_x*y|_i+|G_y*y|_i\bigr),
\end{equation}
where $G_x,G_y$ are Sobel kernels and $\ell_{\mathrm{ce}}(i)$ is the per-pixel cross-entropy. StrongDINOv2Seg uses $\mathcal{L}=\mathcal{L}_{\mathrm{seg}}^{(0)}+\lambda_{\mathrm{aux}}\sum_{k=1}^{3}\mathcal{L}_{\mathrm{seg}}^{(k)}$ with $\lambda_{\mathrm{aux}}=0.4$, where $\mathcal{L}^{(k)}_{\mathrm{seg}}$ is computed on the $k$-th auxiliary output upsampled to input size.

\paragraph{StrongDINOv2Seg.}
The encoder is DINOv2 ViT-B/14~\cite{oquab2023dinov2} (768-d tokens, patch size 14), fully fine-tuned. A $560\times560$ input yields a $40\times40$ token grid. Tokens from transformer blocks 3, 6, 9 and 12 are reshaped to feature maps, projected by $1\times1$ convolutions to 256 channels, and concatenated and fused (FPN-style~\cite{lin2017fpn}) to 512 channels. An ASPP module~\cite{chen2018deeplabv3plus} with rates $\{6,12,18\}$, a $1\times1$ branch and global pooling produces the neck feature. A four-stage decoder with bilinear $\times2$ upsampling and residual convolution blocks (256, 128, 64, 32 channels) is followed by a $1\times1$ classifier and bilinear resampling to input size. Auxiliary $1\times1$ heads attached to the first three decoder stages provide deep supervision~\cite{lee2015dsn} during training. Training uses AdamW (weight decay $10^{-4}$, learning rate $5\times10^{-6}$ for the encoder and $5\times10^{-5}$ for the decoder), cosine annealing to $10^{-7}$ over 10 epochs, batch size 2, gradient-norm clipping at 1.0, and random flips and $90^\circ$ rotations.

\paragraph{From masks to targets.}
Predicted building masks are thresholded, vectorised and cleaned into candidate polygons with a quality class (Section~\ref{sec:data}). The representative point of each candidate becomes a routing target $t\in T$ with weight $w_t=1$; facility points form the set $F$ with weight 1.

\subsection{Graph and Cost Surface}
Let $\Omega_g$ denote the grid cells, $\mathcal{B}\subset\Omega_g$ the cells that are blocked (water, no-data or slope $s(x)>s_{\max}=25^\circ$), and $\mathcal{P}\subset\Omega_g$ the cells intersecting a building footprint of quality class \textsc{Plausible}, \textsc{Review} or \textsc{Unresolved\_Merge}. The planning graph is $G=(V,E)$ with $V=\Omega_g\setminus\mathcal{B}$ and $E$ the 8-neighbour adjacencies with no diagonal corner-cutting through blocked cells. The per-cell cost combines land-cover resistance, slope and a footprint multiplier,
\begin{equation}
c(x)=\bigl[\,c_{\mathrm{lc}}(\ell(x)) + g\bigl(s(x)\bigr)\bigr]\cdot m(x),
\label{eq:cell}
\end{equation}
where $\ell(x)$ is the land-cover class, $c_{\mathrm{lc}}$ a tabulated resistance, $g$ a non-decreasing slope term, and $m(x)$ the footprint multiplier defined below. The weight of an edge $e=(u,v)$ is the mean of its end-cell costs times its step length,
\begin{equation}
w(e)=\tfrac12\bigl(c(u)+c(v)\bigr)\,\lambda(e),\qquad \lambda(e)\in\{10,\;10\sqrt2\}\ \text{m},
\end{equation}
and the cost of a path $\pi=(e_1,\dots,e_n)$ is
\begin{equation}
\Lcost(\pi)=\sum_{j=1}^{n}w(e_j).
\label{eq:lcost}
\end{equation}
Costs are ordinal planning units and not currency.

\paragraph{Local penalty override.}
Let $\mu>1$ be the footprint penalty (here $\mu=4$) and $D(p,r)$ the disc of radius $r$ about a point $p$. For a set of exempted facilities $F_{\mathrm{ex}}\subseteq F$ with radii $\{r_f\}$, the multiplier is
\begin{equation}
m(x)=1+(\mu-1)\,\mathbb{1}[x\in\mathcal{P}]\;\Bigl(1-\mathbb{1}\Bigl[x\in\textstyle\bigcup_{f\in F_{\mathrm{ex}}}D(p_f,r_f)\Bigr]\Bigr).
\label{eq:override}
\end{equation}
The override applies where a facility point lies on the edge of its own footprint, so that the penalty would otherwise make its route detour around its own building. We use it for the health sub-centre (HEALTH\_8734223673) with $r_f=20$~m. Radii of 10, 20 and 30~m were tested; 20~m was the smallest that removed the detour without altering any other route. Crossing statistics are always computed on the original footprints, so the override does not flatter the reported exposure.

\subsection{Shared Network Generation}
\paragraph{Stage A: alternatives.} For each target $t$ we compute a least-cost path $\pi_t^{(1)}$ to the set $R$ of existing road cells (the goal set), where $R$ is restricted to the valid grid. Up to $K-1$ further alternatives are generated by multiplying costs in a corridor of 20~m around previously found paths by factors 3 and 8 and rejecting paths whose overlap with an earlier alternative exceeds 0.8; here $K=3$, giving 319 alternatives for 163 houses.

\paragraph{Stage B: selection with shared-edge credit.} Targets are processed greedily. At each step the candidate that maximises weight over \emph{marginal} cost is chosen, where the marginal cost of $\pi_t^{(k)}$ is
\begin{equation}
\Delta\Lcost\bigl(\pi_t^{(k)}\mid S\bigr)=\sum_{e\in\pi_t^{(k)}\setminus E_S}w(e),
\end{equation}
$E_S$ being the edges of the network selected so far. Edges already built are therefore free, and a route may terminate on another route or on a road, which is how shared access roads emerge. This marginal-cost attachment is the same mechanism as the shortest-path heuristic for Steiner trees~\cite{kou1981steiner,mehlhorn1988steiner} but is applied to an already-penalised, land-cover-weighted graph with an existing road network as a free sub-tree, and it carries no optimality guarantee. After selection, $P$ passes of \emph{leaf re-routing} re-evaluate each leaf route given the rest of the network. The resulting plan is a forest rooted at $R$ with $|T|$ targets and $F$ facilities attached.

Algorithm~\ref{alg:route} summarises the procedure.

\begin{algorithm}[t]
\caption{Cost-surface routing with local penalty overrides}\label{alg:route}
\small
\KwIn{Grid cells $\Omega_g$, land cover $\ell$, slope $s$, footprints $\mathcal{P}$, targets $T$, facilities $F$, road cells $R$, penalty $\mu$, overrides $\{(p_f,r_f)\}_{f\in F_{\mathrm{ex}}}$, slope limit $s_{\max}$, alternatives $K$, passes $P$}
\KwOut{Network edge set $E_S$ and routes $\{\pi_t\}$}
$\mathcal{B}\leftarrow\{x:\ \text{water}\vee\text{no-data}\vee s(x)>s_{\max}\}$\;
\ForEach{$x\in\Omega_g\setminus\mathcal{B}$}{
  $m(x)\leftarrow 1+(\mu-1)\,\mathbb{1}[x\in\mathcal{P}]$\;
  \If{$\exists f\in F_{\mathrm{ex}}:\ \|x-p_f\|\le r_f$}{$m(x)\leftarrow 1$ \tcp*{local override}}
  $c(x)\leftarrow[c_{\mathrm{lc}}(\ell(x))+g(s(x))]\cdot m(x)$\;
}
Build $G=(V,E)$ with $w(e)=\tfrac12(c(u)+c(v))\lambda(e)$\;
\ForEach{$t\in T\cup F$}{
  $A_t\leftarrow\{\pi^{(1)}_t\}$ with $\pi_t^{(1)}=\argmin_{\pi:\,t\to R}\Lcost(\pi)$ \tcp*{Dijkstra/A*}
  \For{$k\leftarrow 2$ \KwTo $K$}{
    inflate $w$ by $\{3,8\}$ inside a 20 m corridor around $A_t$\;
    $\pi\leftarrow$ least-cost path on inflated graph\;
    \lIf{$\max_{\pi'\in A_t}\mathrm{overlap}(\pi,\pi')\le 0.8$}{$A_t\leftarrow A_t\cup\{\pi\}$}
  }
}
$E_S\leftarrow\emptyset$; $U\leftarrow T\cup F$\;
\While{$U\neq\emptyset$}{
  $(t^\star,\pi^\star)\leftarrow\argmax_{t\in U,\,\pi\in A_t}\ w_t\,/\,\Delta\Lcost(\pi\mid E_S)$\;
  $E_S\leftarrow E_S\cup\pi^\star$; $U\leftarrow U\setminus\{t^\star\}$\;
}
\For{$j\leftarrow 1$ \KwTo $P$}{
  \ForEach{leaf route $\pi_t$}{re-route $t$ to the cheapest attachment on $E_S\setminus\pi_t\cup R$; keep if $\Lcost$ decreases}
}
\Return{$E_S,\{\pi_t\}$}
\end{algorithm}

\subsection{Baselines and Evaluation Quantities}
Let $L(\cdot)$ denote the total length of the \emph{unique} edge set of a plan. We compare against (i) \emph{straight-line} links (no terrain); (ii) \emph{independent} least-cost routes, one per house and without sharing; (iii) a \emph{length-only} control (shortest distance under the same hard constraints); and (iv) a \emph{roads-only shared} reference in which houses attach only to existing roads but may share corridors. Savings are
\begin{equation}
\rho=\frac{L_{\mathrm{indep}}-L_{\mathrm{plan}}}{L_{\mathrm{indep}}}\times100\%.
\end{equation}
Exposure to unmodelled buildings is the number of routes crossing a footprint other than the target's own, and the metres of route inside such footprints.

\subsection{Omission Propagation Analysis}
Let $T$ be the full target set and $M\subset T$ an omission set, standing for building false negatives. The planner is re-run on $T\setminus M$, with the footprints of $M$ also removed from the penalty layer, because a missed detection is absent from both the target list and the obstacle layer. We define:
\begin{align}
\kappa &= \frac{|T\setminus M|}{|T|}, \qquad U=|M|,\label{eq:kappa}\\
\delta_L &= \frac{L(T)-L(T\setminus M)}{L(T)},\qquad \bar\ell=\frac{L(T\setminus M)}{|T\setminus M|},\label{eq:delta}\\
E_M &= \text{route length inside footprints of } M \text{ in the plan for } T\setminus M.\label{eq:exposure}
\end{align}
Here $\kappa$ is coverage, $U$ the number of households without a planned route, $\delta_L$ the underestimation of network length, $\bar\ell$ the length per routed target (the quantity that is comparable across scenarios, because absolute length shrinks mechanically when fewer targets are connected), and $E_M$ the exposure to omitted buildings. Omission modes are: random omission of 10\% and 20\% of targets (five draws each); omission of all \textsc{Unresolved\_Merge} candidates ($|M|=62$); omission of the smallest 10\% and 25\% of targets by footprint area (areas $\le46.7$ and $\le113.4$~m$^2$; $|M|=17,41$), a plausible false-negative pattern for small or low-contrast roofs; and random \emph{additions} of 10\% and 20\% synthetic targets as a false-positive analogue. To separate the loss of obstacle information from plain restructuring, every omission run is repeated with $\mu=1$ (no penalty), where the planner ignores footprints and any change in $E_M$ reflects restructuring alone.

\subsection{Hierarchical Multi-Scale Extension}\label{sec:hier}
Level~1 is the village pipeline above. The three coarser tiers reuse its search engine (8-connected Dijkstra with no diagonal corner-cutting) and the shared-network rule of Section~\ref{sec:method}, changing only the nodes, the grid and what counts as an attachment target.

\paragraph{Level 2: Gram Panchayat.}
The 163 house targets are aggregated into hamlet hubs (connected components of houses closer than 150~m; the hub is the member nearest the cluster centroid, its weight the number of houses), which together with the school, sub-centre and two panchayat-office points form the node set. The grid is 10~m, with the footprint penalty $\mu$ and the local override of Eq.~\eqref{eq:override} applied around the amenity points (radius 20~m), and the existing roads are the goal set $R$. Detection omissions act on houses: a fraction of the 163 houses is removed and the hamlets are re-aggregated.

\paragraph{Level 3: mandal.}
Nodes are Gram Panchayat hubs, the Mandal HQ, primary health centres and markets on a 30~m grid. Only collector-grade roads (motorway to tertiary) are attachment targets; lower-class roads merely reduce the cost of travelling along them, and every road cell is traversable (bridges and culverts) at a cost multiplier $0.25$. Water is blocked, tree cover carries an extra multiplier, and slopes above $25^\circ$ are blocked.

\paragraph{Level 4: district.}
A coarse grid (250~m) is built from the same cost recipe with road cells discounted by $0.2$, and the terminals are attached to the root (district HQ) or to one another by the shared rule. The resulting cell tree is contracted into chains between terminals, junctions and leaves (every cell of degree two is eliminated), and each chain is refined at 30~m inside a corridor of 1~km around the coarse path. A chain is cut into segments of at most 4~km; every segment is refined in its own window, with windows aligned to one global 30~m lattice so that neighbouring segments meet in the same cell, and the memory held at any time is the coarse grid plus one window. A window above a cell limit is coarsened rather than failing, and a segment with no fine path keeps its coarse geometry and is flagged. Because a backbone may run along existing roads, we separate new construction from reuse. With $N$ the refined network, $R$ the existing roads and $\mathcal{B}(R,\delta)$ their $\delta$-buffer ($\delta=1.5$ fine cells),
\begin{equation}
L_{\mathrm{new}}=\operatorname{len}\bigl(N\setminus\mathcal{B}(R,\delta)\bigr),\qquad
\rho_{\mathrm{road}}=1-\frac{L_{\mathrm{new}}}{\operatorname{len}(N)},
\label{eq:newconstr}
\end{equation}
and the baseline is the star layout in which every terminal is routed to the root separately.

\paragraph{Penalty sensitivity.}
The exposure of a plan to footprints is the route length lying inside a footprint beyond the first 20~m of each route (the start point sits in or next to its own building),
\begin{equation}
E(\mu)=\sum_{t}\operatorname{len}\Bigl(\bigl(\pi_t\setminus D(p_t,20)\bigr)\cap\mathcal{P}\Bigr),
\label{eq:expo}
\end{equation}
where $D(p,r)$ is the disc of radius $r$ about the start $p_t$ of route $\pi_t$ and $\mathcal{P}$ the footprints. We evaluate $\mu\in\{1,2,4,8,20,100\}$ at Level~2 and report $E(\mu)$ and the network length relative to $\mu=1$.

% -----------------------------------------------------------------------------
\section{Experimental Results and Discussion}\label{sec:results}

\subsection{Segmentation Performance}\label{sec:segres}
Table~\ref{tab:seg} reports the segmentation evaluation. StrongDINOv2Seg was validated on the 1{,}669-image LoveDA validation split (Rural and Urban domains, seven scored classes, ID~0 ignored) at 560~px, and reached 53.8\% mIoU, 68.9\% Dice, 71.3\% pixel accuracy, 72.4\% macro precision and 66.9\% macro recall at its best epoch (4 of 10; pixel accuracy ranged from 69.0\% to 72.2\% over the ten epochs). These are \emph{source-domain} figures; they are not accuracy on Damarakunta, for which no pixel reference exists. The epoch-4 per-class results for \emph{Building} are 54.0\% IoU, 62.1\% precision and 80.6\% recall: roughly one building pixel in five is missed and more than one in three predicted building pixels is spurious in the source domain, which is the pixel-level counterpart of the omission scenarios analysed below. 

The lower block of Table~\ref{tab:seg} gives the only village-level evidence: agreement of the 165 in-boundary candidates with Open Buildings. Overall agreement is 71.5\% (118/165) but varies from 43.8\% for \textsc{Review} to 85.5\% for \textsc{Unresolved\_Merge}. The high agreement of merged blobs is consistent with merges being large, conspicuous roofs; it says nothing about whether the merged polygon contains one house or several, which is exactly what matters for per-household routing. In the other direction Open Buildings lists 1{,}537 polygons inside the boundary, of which 1{,}186 have no candidate, so the candidate set must not be read as a complete house count. Part of that gap is expected (sheds and farm structures), but the cause was not established; the 163 targets are therefore a conservative planning set.

\begin{table}[t]
\centering
\caption{Segmentation evaluation. Top: LoveDA validation (source domain). Bottom: agreement of Damarakunta candidates with Open Buildings v3 (a proxy; not accuracy). Precision and recall are macro-averages over the seven scored classes; Dice for the Building class is its F1.}
\label{tab:seg}
\footnotesize
\setlength{\tabcolsep}{2.5pt}
\begin{tabular}{@{}lccccc@{}}
\toprule
\multicolumn{6}{@{}l}{\textit{LoveDA validation split (\%)}}\\
Model & mIoU & Dice & Pixel Acc. & Prec. & Rec.\\
\midrule
StrongDINOv2Seg & 53.8 & 68.9 & 71.3 & 72.4 & 66.9\\
\quad Building class & 54.0 & 70.1 & --- & 62.1 & 80.6\\
\midrule
\multicolumn{6}{@{}l}{\textit{Damarakunta candidates vs.\ Open Buildings v3}}\\
Class & Cand. & Agree & Disagree & \multicolumn{2}{c}{Agreement}\\
\midrule
All & 165 & 118 & 47 & \multicolumn{2}{c}{71.5\%}\\
\textsc{Plausible} & 71 & 51 & 20 & \multicolumn{2}{c}{71.8\%}\\
\textsc{Review} & 32 & 14 & 18 & \multicolumn{2}{c}{43.8\%}\\
\textsc{Unresolved\_Merge} & 62 & 53 & 9 & \multicolumn{2}{c}{85.5\%}\\
\bottomrule
\end{tabular}
\end{table}

\subsection{Routing Efficiency and Detour Minimisation}
Table~\ref{tab:routing} compares the planning strategies for the 163 routed houses. Independent least-cost routes total 10{,}045~m, but because adjacent houses re-trace the same corridors the unique built length is 6{,}753~m. The two-stage plan reduces this to 5{,}069~m, a saving of 1{,}684~m (24.9\%), and all 163 houses and all facilities are connected with no hard-constraint violation, no cycles and no route crossing without a shared cell. A length-only control, which minimises distance under the same constraints, is \emph{longer} in unique length than independent least-cost routing ($+3.4$\%) and crosses other footprints on 64 routes, so the cost surface, rather than shortness alone, produces the gain. The roads-only shared reference (5{,}229~m) is 160~m longer than the two-stage plan, showing that routes ending on other routes, not only on roads, contribute a modest additional saving; the ablation over $K\in\{1,3\}$ and $P\in\{0,2\}$ shows that extra alternatives change nothing in this village ($K=1$ and $K=3$ give identical lengths) and that leaf re-routing contributes only 10.3~m (0.2\%), so the saving comes mainly from sharing itself.

\begin{table*}[t]
\centering
\caption{Routing efficiency for 163 houses (building penalty $\mu=4$ unless stated). Total length counts every route; unique length counts shared edges once. Crossings count routes passing through another house's footprint.}
\label{tab:routing}
\small
\begin{tabular}{@{}lrrrrr@{}}
\toprule
Strategy & Total (m) & Unique (m) & $\rho$ vs.\ indep.\ (\%) & Routes crossing & Length in footprints (m)\\
\midrule
Straight line (2 targets infeasible) & 8{,}092.3 & --- & --- & 60 & 1{,}453.2\\
Independent least-cost routes & 10{,}045.4 & 6{,}753.1 & 0.0 & 46 & 328.3\\
Length-only control & 9{,}166.3 & 6{,}979.8 & $-3.4$ & 64 & 1{,}615.8\\
Roads-only shared reference & 5{,}229.3 & 5{,}229.3 & 22.6 & 29 & 216.1\\
\textbf{Two-stage plan (proposed)} & \textbf{5{,}069.0} & \textbf{5{,}069.0} & \textbf{24.9} & \textbf{28} & \textbf{209.1}\\
\midrule
\multicolumn{6}{@{}l}{\textit{Control: no footprint penalty ($\mu=1$)}}\\
Independent least-cost routes & 9{,}595.9 & 6{,}921.3 & 0.0 & 63 & 1{,}423.1\\
Two-stage plan & 4{,}705.5 & 4{,}705.4 & 32.0 & 56 & 660.1\\
\midrule
\multicolumn{6}{@{}l}{\textit{Hard building avoidance via 7.1~m buffer, two-stage plan}}\\
Penalty $\times20$ & --- & 6{,}238.1 & 18.5 & 19 & 147.9\\
Penalty $\times100$ & --- & 6{,}190.3 & 18.9 & 19 & 144.8\\
\bottomrule
\end{tabular}
\end{table*}

The footprint penalty exposes a trade-off between corridor length and obstacle avoidance. Without any penalty ($\mu=1$) the plan is shortest (4{,}705~m), but 56 routes pass through other houses (660~m inside footprints). With $\mu=4$ the plan is 363~m (7.7\%) longer but crossings halve to 28 routes and 209~m. Buffered penalties push crossings down to 19 routes and 145~m but cost a further 1.1--1.2~km of corridor and erode the saving to 18.5\%. Hard blocking of footprint cells was not feasible because every house point lies inside its own blocked cell. Remaining crossings under the buffer scenario are mostly confined to the first 20~m of the route, i.e.\ the start of the route leaving its own house.

\paragraph{Healthcare node and the local penalty override.}
The health sub-centre point lies 1.8~m from a 1{,}194~m$^2$ unresolved-merge footprint. Without an override, the $\times4$ penalty makes its route detour around that footprint. Table~\ref{tab:override} shows that lifting the penalty within 20~m shortens the access route from 73.3~m to 67.4~m and lowers its cost from 357.7 to 260.1 units ($-27.3$\%), while 0 of 163 house routes change and the combined unique length of houses and facilities falls from 5{,}324.2~m to 5{,}318.3~m. The effect is local by construction and small in metres, but the cost reduction is large, and because the override changes nothing else it can be audited in isolation. The route still ends on the road: facilities are planned before houses, so they cannot attach to house routes. We stress that the override is a manual modelling decision, not a finding about access.

\begin{table}[t]
\centering
\caption{Effect of the local penalty override at the health sub-centre ($r=20$~m).}
\label{tab:override}
\small
\begin{tabular}{@{}lrrr@{}}
\toprule
 & Without & With & Change\\
\midrule
Sub-centre route length (m) & 73.3 & 67.4 & $-8.0\%$\\
Sub-centre route cost (units) & 357.7 & 260.1 & $-27.3\%$\\
House routes changed & --- & 0 of 163 & ---\\
Combined unique length (m) & 5{,}324.2 & 5{,}318.3 & $-0.1\%$\\
\bottomrule
\end{tabular}
\end{table}

\subsection{Impact of Building Omissions on Network Coverage}
Table~\ref{tab:omission} summarises the omission analysis against the 5{,}069~m baseline. Three observations follow.

\emph{First, the plan becomes shorter, not longer.} Missing buildings make the corridor budget \emph{look} smaller: 10\% random omission (16 houses) understates length by 7.6\%, 20\% by 16.9\%, omission of the smallest quarter by 25.1\% and omission of all merged blobs by 29.4\%. A planner reading only total length would therefore underfund the scheme.

\emph{Second, the households that vanish from the plan are the ones that lose access.} $U$ households receive no route: 16 and 33 for random 10\% and 20\% omission, 41 for the smallest quarter, and 62 (38.0\%) for merged blobs. The omitted sets are not spatially random in the area-based modes, and merged blobs correspond to dense multi-roof structures; if these coincide with the dense core of a hamlet, the loss of coverage is concentrated where need is greatest. We did not measure that spatial concentration, so it is stated as a risk, not a result.

\emph{Third, length per routed target is the comparable quantity, and it moves in a direction that hides the problem.} Per-target length is stable under random omission (31.9~m and 32.4~m against 31.1~m) but rises by 14\% when merged blobs are removed (35.4~m), because the remaining targets are the peripheral ones. Conversely, random additions increase total length by 8.6\% and 15.8\% and lower per-target length, since dense additions share existing corridors. Savings $\rho$ against independent routing stay within about 1.7 percentage points of the baseline in every scenario, so the benefit of sharing is robust even when the absolute length is not.

\begin{table*}[t]
\centering
\caption{Impact of simulated building omissions (and additions) on the shared network. Baseline: $L=5{,}069.0$~m, 163 routed. Random scenarios are averages of five draws. $U$: households without a planned route (relative to the 163 baseline targets). $\delta_L$: length change relative to baseline (negative: shorter). $\bar\ell$: length per routed target.}
\label{tab:omission}
\small
\begin{tabular}{@{}lrrrrrrr@{}}
\toprule
Scenario & Routed & $U$ & $\kappa$ (\%) & $L$ (m) & $\delta_L$ (\%) & $\bar\ell$ (m) & $\rho$ vs.\ indep.\ (\%)\\
\midrule
Baseline & 163 & 0 & 100.0 & 5{,}069.0 & 0.0 & 31.1 & 24.9\\
Omit 10\% (random) & 147.0 & 16 & 90.2 & 4{,}684.4 & $-7.6$ & 31.9 & 25.1\\
Omit 20\% (random) & 130.0 & 33 & 79.8 & 4{,}211.5 & $-16.9$ & 32.4 & 25.0\\
Omit smallest 10\% & 146 & 17 & 89.6 & 4{,}648.7 & $-8.3$ & 31.8 & 24.9\\
Omit smallest 25\% & 122 & 41 & 74.8 & 3{,}798.3 & $-25.1$ & 31.1 & 25.6\\
Omit merged blobs & 101 & 62 & 62.0 & 3{,}579.0 & $-29.4$ & 35.4 & 26.6\\
\midrule
Add 10\% (random) & 177.4 & --- & --- & 5{,}502.4 & $+8.6$ & 31.0 & 25.6\\
Add 20\% (random) & 190.6 & --- & --- & 5{,}868.3 & $+15.8$ & 30.8 & 25.4\\
\bottomrule
\end{tabular}
\end{table*}

\paragraph{Does the planner's information loss matter?}
Removing footprints from the penalty layer also removes information that steers routes around houses. On an earlier baseline run of the same planner (5{,}137.7~m), length inside omitted footprints rose from 52.6~m to 178.0~m when merged blobs were omitted, and routes crossing an omitted footprint rose from 4 to 17. For the smallest-quarter scenario the increase was modest (25.9 to 33.3~m), and for the smallest 10\% it fell (14.9 to 3.6~m), so the effect is not monotone. With the penalty disabled, the same comparisons fall (148.3 to 113.5~m for merged blobs), which indicates that the primary-run increase is attributable to lost obstacle information, whereas the control isolates restructuring. These are single deterministic scenarios per mode; they indicate direction and rough magnitude, not effect sizes with uncertainty. They also show that omission errors act twice: they drop households from the plan and let corridors run through the dwellings that were omitted.

\paragraph{Why IoU is not enough.}
The segmentation metric of Table~\ref{tab:seg} is area-weighted over classes, whereas the planner counts \emph{instances}. A model that misses the smallest quarter of buildings by count can lose only a small share of building pixels, yet it removes a quarter of the households from the plan. Conversely a model that merges adjacent roofs scores a high pixel IoU but, as the merged-blob scenario shows, loses 38\% of the targets in a single stroke. We therefore recommend that detection studies intended for infrastructure use report instance-level recall stratified by footprint size and by merge status, alongside mIoU and Dice.

\subsection{Gram Panchayat Level and the Footprint-Penalty Sweep}\label{sec:gp}
Level~2 is the second real-data tier: the Damarakunta Gram Panchayat with the repository's TGRAC roads, the five facilities, the v8 building footprints, the 163 house targets, WorldCover and the Copernicus DEM. Its 18 hamlet hubs and 5 facilities connect with 1{,}743.7~m of shared network against 1{,}782.3~m for independent routing (a 2.2\% saving; Table~\ref{tab:gp}). The small saving is expected: 12 of the 18 hubs (median 60~m) already lie within 100~m of a road, so there is little corridor left to share. Hamlet-hub lengths are not comparable with the 5{,}069~m of Level~1, which routes 163 individual houses. Lifting the footprint penalty within 20~m of an amenity (Eq.~\eqref{eq:override}) shortens two facility routes by 22.4 and 20.0~m (42.4~m in total), lowers the ordinal cost of all five facility routes by 33--67\%, and changes none of the other 18 hub routes. Removing 10\% and 20\% of the houses at random (three seeds) shortens the plan by 5.9\% (range 2.3--8.5\%) and 11.1\% (7.3--14.0\%) and leaves 16 and 33 houses unserved, the same pattern as at village level.

\begin{table}[t]
\centering
\caption{Level 2: Damarakunta Gram Panchayat, real inputs (10~m grid, $\mu=4$). Omission rows are means of three seeds (11, 12, 13) with the range in brackets; units are houses.}
\label{tab:gp}
\footnotesize
\setlength{\tabcolsep}{3pt}
\begin{tabular}{@{}lr@{}}
\toprule
Quantity & Value\\
\midrule
Nodes (hamlet hubs + facilities) & 18 + 5\\
Houses represented by the hubs & 163\\
Shared unique length (m) & 1{,}743.7\\
Independent unique length (m) & 1{,}782.3\\
Saving vs.\ independent (\%) & 2.2\\
Override detour reduction (m) & 42.4\\
Other routes changed by the override & 0\\
Omit 10\% of houses: $\delta_L$ (\%), unserved & 5.9 [2.3--8.5], 16\\
Omit 20\% of houses: $\delta_L$ (\%), unserved & 11.1 [7.3--14.0], 33\\
\bottomrule
\end{tabular}
\end{table}

The sweep of Table~\ref{tab:sweep} shows that the penalty has a knee. Raising $\mu$ from 1 to 4 removes 89.0\% of the route length inside footprints (109.1 to 12.0~m) and cuts the routes crossing a footprint from 6 to 3 for 5.2\% more network length. Going on to $\mu=8$ buys 4.4 more points of exposure reduction for 2.9 more points of length, and $\mu\ge 20$ changes the exposure no further while the length keeps growing. We therefore keep $\mu=4$, the value used at Level~1, as the primary penalty and report the sweep as a sensitivity. At Level~1 the same direction appears: removing the penalty ($\mu=1$) gives the shortest plan (4{,}705~m) but 56 routes through other houses, whereas $\mu=4$ halves the crossings (Table~\ref{tab:routing}).

\begin{table}[t]
\centering
\caption{Footprint-penalty sensitivity at Level 2 (real inputs). $E(\mu)$ is the route length inside footprints (Eq.~\eqref{eq:expo}); changes are relative to $\mu=1$.}
\label{tab:sweep}
\footnotesize
\setlength{\tabcolsep}{2.5pt}
\begin{tabular}{@{}rrrrrr@{}}
\toprule
$\mu$ & Length (m) & $\Delta L$ (\%) & Routes crossing & $E(\mu)$ (m) & $E$ reduction (\%)\\
\midrule
1 & 1{,}657.1 & 0.0 & 6 & 109.1 & 0.0\\
2 & 1{,}692.0 & 2.1 & 4 & 45.6 & 58.2\\
\textbf{4} & \textbf{1{,}743.7} & \textbf{5.2} & \textbf{3} & \textbf{12.0} & \textbf{89.0}\\
8 & 1{,}791.2 & 8.1 & 2 & 7.2 & 93.4\\
20 & 1{,}813.7 & 9.5 & 2 & 7.2 & 93.4\\
100 & 1{,}842.0 & 11.2 & 2 & 7.2 & 93.4\\
\bottomrule
\end{tabular}
\end{table}

\subsection{Mandal and District Levels: A Scalability Demonstration}\label{sec:scale}
\emph{The mandal and district runs are a computational and architectural demonstration, not planning findings.} Their terminals are proxies (Section~\ref{sec:opendata}): village points as Gram Panchayat hubs, the interior point of a polygon as the Mandal HQ and as every mandal terminal, and an interior point of the Siddipet Urban polygon as the district HQ. Their only purpose is to show that the cost recipe, the shared rule and the coarse-to-fine machinery run on real rasters and real road layers at the scale of a mandal (116~km$^2$) and a district (3{,}722~km$^2$) without memory growth.

At the mandal level, 14 village points and the proxy Mandal HQ attach to collector-grade roads with 5{,}913~m of new link; the villages are scattered, so shared and independent plans coincide (Table~\ref{tab:scale}). At the district level, 22 proxy terminals connect to the proxy district HQ with a backbone of 358.7~km, of which 344.1~km (95.9\%) already exists as road and only 14.7~km is new construction; the star layout would need 19.8~km of new construction, so sharing saves 25.9\% of it. The total-length saving over the star (22.2\%) would overstate the benefit, which is why Eq.~\eqref{eq:newconstr} separates reuse from construction. The district run used 112{,}847 coarse cells and 105 fine windows (all refined at 30~m, none degraded or falling back to the coarse path) and took about 47~s on one CPU core. Removing 10\% and 20\% of the units at random changes the district backbone by 7.6\% (4.8--10.4\%) and 11.5\% (4.2--16.3\%), and the Markook plan by 2.1\% and 25.5\% (the latter ranging 5.4--63.2\% because only three of the 14 village units are removed); these omission rows only exercise the sensitivity code on proxy nodes.

\begin{table*}[t]
\centering
\caption{Levels 3 and 4: scalability demonstration on \emph{proxy} nodes (real rasters, Overture roads). Not planning findings. Omission: mean $\delta_L$ (\%) of three seeds [range] for 10\% / 20\% of units removed.}
\label{tab:scale}
\footnotesize
\begin{tabular}{@{}lllrrrrrl@{}}
\toprule
Level & Nodes (proxy) & Grid & Shared (km) & On existing (km) & New (km) & Star, new (km) & Sharing saves new (\%) & Omission $\delta_L$\\
\midrule
3 Mandal (Markook, 116~km$^2$) & 14 villages + 1 HQ & 30~m & 5.9 & -- & 5.9 & 5.9 & 0.0 & 2.1 [1.2--3.3] / 25.5 [5.4--63.2]\\
4 District (Siddipet, 3{,}722~km$^2$) & 22 mandal terminals & 250~m / 30~m & 358.7 & 344.1 & 14.7 & 19.8 & 25.9 & 7.6 [4.8--10.4] / 11.5 [4.2--16.3]\\
\bottomrule
\end{tabular}
\end{table*}

\subsection{Implications for Rural Infrastructure Planners}
For PMGSY-style schemes, three practical points emerge. (i)~\emph{Treat building inventories as uncertain planning inputs.} A network sized from a candidate layer should carry a sensitivity band; in this village a 10--20\% omission corresponds to a 7.6--16.9\% length understatement and 16--33 unserved households. (ii)~\emph{Prefer shared corridors.} Sharing saves roughly a quarter of the unique length against house-by-house links and this saving is stable under omissions, which makes it a robust planning argument even where inventories are poor. (iii)~\emph{Use local overrides sparingly and visibly.} A facility whose point lies on a footprint edge can otherwise force a detour in a high-priority route; a bounded override fixes this without changing other routes, and its effect can be reported separately. All outputs are \emph{planning alternatives}: costs are ordinal, no field survey, land ownership, legal access, road width or construction cost is modelled, and the plans are not substitutes for engineering survey.

\subsection{Threats to Validity}
The study covers one village and the omission scenarios are synthetic, not measured segmentation false-negative rates. The candidate layer has no surveyed reference. The LoveDA figures are source-domain results from the epoch-4 checkpoint (metadata recovered from the training notebook), and the checkpoint's hash was not logged with the village candidates. No model comparison on Damarakunta is claimed. The planner is greedy and not globally optimal; shared-edge credit counts only identical grid edges. Slope derives from a 30~m DEM. About 45\% of the reference road length lies outside the village boundary; one house (C0002) attaches to a road component that is not the main component inside the clipped domain, which we regard as a boundary-truncation artefact rather than a planner defect. The imagery licence for derived products has not been cleared. For the hierarchical tiers, the Gram Panchayat result routes hamlet hubs on a real land-cover and DEM surface, the mandal and district tiers rest on proxy nodes and OSM-derived roads of unknown completeness, the geoBoundaries release may not match the current mandal set, and Levels~2--4 use a different DEM (Copernicus GLO-30) from Level~1 (SRTM), a difference we quantify but do not remove. Omission scenarios at every tier are synthetic. Costs are ordinal units; no rupee construction cost is modelled yet.

% -----------------------------------------------------------------------------
\section{Conclusion and Future Work}\label{sec:conclusion}
We presented an end-to-end pipeline from satellite imagery to shared rural access roads, comprising a fine-tuned DINOv2 segmenter, a penalised multi-criteria cost surface with a local override and a two-stage generate-and-select planner, and used it to quantify how building omissions propagate into network length, topology and household coverage in Damarakunta. The shared plan connects all 163 houses with 24.9\% less unique corridor than independent routing; the sub-centre override removes a detour (cost $-27.3$\%) without changing any other route. Omissions shorten the plan, drop households from it, and make corridors cross the buildings that were missed; in the worst scenario length is understated by 29.4\% and 38.0\% of households are unserved, while pixel metrics would suggest a competent model. Planners should read detection quality through instance recall, per-target length and coverage. We then lifted the same logic to Gram Panchayat, mandal and district tiers: at Gram Panchayat level on real inputs a footprint penalty of 4 is the knee of the length--exposure trade-off (89.0\% less exposure for 5.2\% more length), and the coarse-to-fine district run shows that the approach scales to 3{,}722~km$^2$ while exposing that 95.9\% of an open-data backbone is existing road. The mandal and district tiers use proxy nodes and are a scalability demonstration only.

\paragraph{Limitations.} Besides the threats above, seasonal occlusion (canopy, shadows, crop cover) is not modelled and affects both detection and land-cover cost; DEM resolution limits the terrain term; the footprint penalty is a tunable multiplier rather than a measured resistance; and a rigorous evaluation requires a surveyed reference for at least one village.

\paragraph{Future work.} (i)~Evaluate the segmenter, and the CNN baselines from earlier prototyping, under one protocol on a hand-labelled Damarakunta subset and measure real false-negative patterns instead of synthetic ones. (ii)~Propagate detection probabilities through the planner (Monte-Carlo or expected-length formulations) to produce uncertainty bands. (iii)~Add multi-modal inputs, namely SAR for cloud and season robustness, a higher-resolution DEM, and mobile-phone or field-survey priors. (iv)~Replace the greedy selector by a Steiner-tree solver on the alternative graph, and attach facilities to house routes. (v)~Test transfer to further villages and to hill and tribal terrain where PMGSY thresholds differ. (vi)~Replace the proxy hubs by verified Gram Panchayat polygons with house counts and verified Mandal HQ, PHC and market coordinates, so that Levels~3 and 4 become planning results, and replace ordinal cost by a rupee model based on a State Schedule of Rates.

% -----------------------------------------------------------------------------
\section*{Data and Code Availability}
Planning outputs, validation manifests and scripts are in the project repository; the imagery mosaic is not redistributable without licence clearance. Every number in Section~\ref{sec:results} is mapped to its input hashes, command, seed and output file in \texttt{docs/results/PROVENANCE\_AND\_PROOFS.md}, which is generated and checked by \texttt{python -m src.evaluation.provenance\_audit}; the hierarchical results are in \texttt{docs/results/hierarchy\_real\_run/}.

\bibliographystyle{IEEEtran}
\bibliography{references}

\end{document}
